% Folder 74, batch 04 — archivists' pages 61–80.
% Pass: Opus 5.5 (claude-opus-5-5), 2026-10-03 — first pass, unchecked against the pages by a human.
%
% Pages 67 (typed Bourbaki seminar programme, no hand) and 71 (blank coloured
% separator sheet) carry no mathematics and are skipped. Page 77 is a crowded
% landscape sheet of diagrams: its list of sets is transcribed, the diagrams
% are described. The connective prose is often a fast hand; the formulas
% carry better than the words between them, and the apparatus says so.
\documentclass[11pt,a4paper]{article}
\input{../preamble/grothendieck}

\folder{74}
\batch{4}
\pages{61}{80}
\foldertitle{Complexes cubiques : notes manuscrites (s.d.).}
\dating{[à partir de 1976]}
\watermark{Édition de démonstration}

\begin{document}

\section{Bitriangles, triades et hexagones}

\page{61}
\note{la page continue un argument commencé avant ce lot : les $X_\alpha$, les fibres $I_\alpha$ de $I \to \Lambda$ et les triades « de type I, II, III » y sont supposés connus.}

Voici un hexagone \add{(\struck{\ill{}} \ill{}) dont les \uncertain{milieux} (deux à deux) \uncertain{non liés}} \struck{de triades \ill{}} contenu dans aucun des $X_\alpha$ [\uncertain{i.e. rencontrant} $(i,j,k)$, $(p,q,r)$, $(s,t,u)$, les trois fibres de $I \to \Lambda = \{1,2,3\}$]

\drawing{un hexagone dont les six sommets sont des triades, écrites comme triples de paires d'éléments pris dans les trois fibres ; en haut, marqué II : $(ip)(\cdot\, j)(uk)$, l'élément du milieu illisible ; à gauche, en haut, marqué I : $(ir)(qs)(qt)$ ; à droite, en haut, marqué I : $(iq, rs, rt)$ ; à gauche, en bas, marqué I : $(\cdot, km, jr)$, la première paire surchargée ; à droite, en bas, marqué I : $(ru)(kq)(\cdot\, q)$ ; en bas, marqué II : $(pu)(ti)(si)$. La lecture des paires est en grande partie douteuse.}

Tout triangle de triades \add{bi-liées} \struck{est} formé, \uncertain{soit} de triades d'un même $X_\alpha$ (de réunion $X_\alpha$), soit de triades de type I et II \add{(2 de type I, et 1 de type II)}, \add{soit des} triades de type III. \struck{Donc tout} L'intersection du bitriangle correspondant avec deux des $X_\alpha$ est un doublet, avec le troisième c'est une étoile de triangles : deux bandes

\drawing{deux triangles ayant le sommet $a$ en commun.}

\struck{Cela \ill{}} \struck{[Pour une \add{étoile} \ill{} donnée dans $X$ il y a exactement trois bitriangles de $S$ coupant $X$ suivant $E$. Pour un doublet donné dans $X$, il y a exactement $2\times 3$} \struck{$(3 \times \frac{3!}{2}) = 54$]}
\note{le passage entre crochets est barré de trois longues diagonales.}

\drawing{figure à neuf sommets $i, j?, k, p, q, r, s, t, u$ : un rectangle $i\,p\,q\,k$ coupé par des traits pleins et pointillés, l'arête $ip$ marquée $(a)$, et une colonne verticale $s, t, u$ reliée en $u$ aux sommets $p$, $q$, $r$.}

\page{62}
Les bitriangles sont en corr. biunivoque avec l'ens. des \add{couples $(a,u)$ d'un} \struck{\ill{}} des sommets de $\Delta$, \add{et d'}une \struck{\ill{} \ill{} \ill{}} \struck{\ill{} \ill{} \ill{}} \uncertain{bijection} $u :$ \struck{$\Pi$}$(X_\alpha) = I_\alpha$
\note{le symbole biffé devant $(X_\alpha)$ ressemble à un $\Pi$ ou à un II.}
où $X_\alpha \ni a$ \uncertain{et} $I(X_\alpha) = I_\alpha$ est l'ens. des trois \uncertain{arêtes} \uncertain{associées}. À un couple $(a,u)$ correspond le bitriangle formé de $a$, des 4 sommets \ill{} liés à $a$ dans $X_\alpha$, et des 4 sommets de $S$ qui sont dans l'\uncertain{arête} choisie $u$, et qui sont \uncertain{non} liés à $a$. Pour un $a \in X = X_\alpha$ donné, il y a donc exactement 3 bitriangles $Y$ qui coupent $X$ suivant la réunion des deux triangles de $X$ \uncertain{passant par} $a$ ; et pour un doublet de $X$ de sommets $\{a,b\}$ de $X$ ($a$ et $b$ non liés) il existe exactement \struck{\ill{}} \struck{3} bitriangles $Y$ coupant $X$ en $\{a,b\}$. [\uncertain{Total} donc :

\begin{itemize}
  \item \struck{les trois}
  \item le bitriangle $X$ — 1
  \item les 2 bitriangles $X'$, $X''$ disjoints de $X$ — 2
  \item les $3 \times 9$ bitriangles $Y$ coupant $X$ suivant une 2-étoile de triangles — 27
  \item les $3 \times$ \struck{$10$} \add{$18$} bitriangles $Y$ coupant $X$ suivant un doublet — 54
  \item les \struck{\ill{}} \add{$6 \times 6$} bitriangles $Y$ coupant $X$ suivant un triangle — 36
\end{itemize}

\struck{NB il y a} — total \struck{\ill{}} 120

On avait \struck{\ill{}} négligé de détailler le cas des $Y$ \uncertain{coupés} par des triades de type III, plus simple.

\page{63}
On trouve qu'ils correspondent au cas de \struck{\ill{}} la figure ci-contre, correspondant à un \uncertain{syst.} transitif d'isom. entre les $I_\alpha$ (définis par les triangles intersection $Y \cap X_\alpha$). Pour \uncertain{tt} triangle donné de $X_\alpha$ (correspondant à une bijection $I_\beta \to I_\gamma$) il y en a 6.

\drawing{un prisme vu en perspective, coupé en trois étages par des triangles horizontaux, les arêtes verticales reliant les étages.}

\uncertain{Étudier} plus généralement ?
\begin{itemize}
  \item a) Position relative de 2 $\mathcal{T}_i$ \ill{} dans un $\mathcal{T}_2$
  \item b) Position relative d'une \uncertain{base} et d'un $\mathcal{T}_1$.
\end{itemize}
\note{la lettre transcrite $\mathcal{T}$ est une majuscule cursive en forme de Z barré, de lecture douteuse ; elle revient aux pages 70 et 75.}

(20) \uncertain{(commentaire)} Soit $\tau = \{a,b,c\}$ une triade du $\mathcal{T}_2$ $\Delta = (S,A)$. Alors \struck{\ill{}} il y a exactement un bitriangle contenant $\tau$, qui est aussi formé des él. de $\tau$ et des 6 él. de $S \setminus \tau$ qui sont liés à \add{exactement} deux él. de $\tau$ \struck{(ils sont liés aux \ill{} \ill{}, deux tels éléments)}. Alors \struck{$X$} $\setminus$ \struck{$\tau$} $= Y$ est un \struck{hexagone} \uncertain{distingué}, \ill{} un hexagone. On trouve ainsi une bijection de l'ens. des triades sur l'ens. des hexagones.

\page{64}
La bijection inverse est celle qui associe à \uncertain{tt} hexagone $Y$ l'ens.\ $\tau$ des $s \in S \setminus Y$ qui sont liés \ill{} \ill{} un élément de $Y$ (ils sont liés à exactement 4 éléments \add{\uncertain{opposés deux à deux}} de l'hexagone, 2 sur chacun de ses triangles inscrits, et ils correspondent biunivoquement aux trois diagonales de l'hexagone).

Autre procédé (qui revient au précédent, en remplaçant la triade $\tau$ par son \uncertain{opposée} $\bar\tau$) : on prend l'ens. des $s \in S \setminus \bar\tau$ qui sont non liés à tous les él. de $\bar\tau$, \struck{\ill{}} c'est \struck{\ill{}} un hexagone, \struck{et \ill{} \ill{}} \uncertain{que précédemment} \ill{} par $\bar\tau$ est l'ens. des él. de $S$ liés à tous les él. de $\tau$).

\uncertain{On examine} pour une triade donnée $\tau$, et $x \in S - \tau$, regardons \struck{\ill{}} card $\mathrm{Ét}(x) \cap \tau$ \add{\uncertain{ens. des $s \in \tau$ liés à $x$}}, soit $\nu_\tau(x)$.
\note{le symbole lu $\mathrm{Ét}$ est douteux.}
\uncertain{Il y a} :
\begin{itemize}
  \item 6 él. tels que $\nu_\tau(x) = 0$, ils forment un hexagone $h(\tau)$ ;
  \item 9 él. tels que $\nu_\tau(x) = 1$, ils forment un bitriangle $\mathrm{bit}(\tau)$ ;
  \item 6 él. tels que $\nu_\tau(x) = 2$, ils forment un hexagone $\bar h(\tau)$ (dit « opposé » à $h$) ;
  \item 3 él. tels que $\nu_\tau(x) = 3$, ils forment une triade $\bar\tau$ (dite « opposée » à $\tau$).
\end{itemize}

\page{65}
On a
\[
  \mathrm{bit}(\tau) = \mathrm{bit}(\bar\tau) = \mathrm{bit}(\tau') = \mathrm{bit}(\tau'')
  = \mathrm{bit}(\bar\tau') = \mathrm{bit}(\bar\tau'')
\]
\note{le troisième terme porte sous $\tau'$ un indice douteux, peut-être $6$.}
où $\tau', \tau''$ sont les triades \uncertain{composantes} de l'hexagone $h(\tau)$, de sorte que $\bar\tau', \bar\tau''$ sont les triades composantes de l'hexagone $h(\bar\tau) = \bar h(\tau)$.

\marginal{détailler tout ça dans \uncertain{énoncés} en forme….}

\page{66}
\drawing{un hexagone à sommet central $b$ ; sommets étiquetés, en partie surchargés : en haut à gauche $b' = A_4$, en haut à droite $\cdot = A_3$, à gauche $a'$ (précédé de « 2: »), à droite $\cdot = A_2$, en bas à gauche $c = A_0$, en bas à droite $a = A_1$ ; les arêtes $a'b'$, $A_0A_1$ et $A_1A_2$ sont repassées en gras. À côté, deux petits schémas croisant des lignes entre $\xi_1, \xi_2, \xi_3$ et $\xi_1^\wedge, \xi_2^\wedge, \xi_3^\wedge$ (lecture douteuse).}

$\tau' \in \tilde E_c$ \uncertain{lié} à \ill{} et à \ill{}

$27 \cdot 10 \cdot 8 \cdot 4 \cdot 1 \cdot 2$
\note{le 7 et le 10 sont repassés.}

$2^7 \cdot 3^3 \cdot 5$ = Nb de hexagones \uncertain{épinglés}

$2^6 \cdot 3 \cdot 5$ = Nb de hexagones = Nb de bitriades

\uncertain{Elles} sont conjuguées

La donnée d'un $A_4$ (\struck{$\Leftrightarrow$} d'un $A_5$) est équivalente à celle d'un bitriangle épinglé dans $S$

\struck{$(a,b,c)$} \quad \struck{$(x,y,z)$} \quad \struck{\ill{}}

\[
  A_4 \Leftrightarrow A_5 \Leftrightarrow P_6 \Leftrightarrow \text{bitriangle épinglé}
\]
bitriades \uncertain{sont} conjuguées
\struck{Nb de bitriades épinglées}

\struck{\ill{}} \struck{Nb de bitriades} \struck{Nb de} card bitriad $= \left(\frac{12}{72}\right) \frac{1}{6}$ Nb des bitriangles

\struck{\ill{}} $27 \cdot 16 \cdot 10 = 6 \cdot 720$ ✓
\struck{$27 \cdot 16 \cdot 8 = 3 \cdot 720$} = Nb des triades épinglées

NB \quad $45 \times 16 = 720$ = Nb de bitriangles \uncertain{munis} d'un triangle \struck{= Nb de triades} = Nb de triades

Nb de bitriangles $= 120$

\section{Pavés et triades}

\page{68}
$\Delta = (S,A,T)$ graphe \uncertain{cubique}

$X := X_0 \subset S$ pavé (anciennement : bitriangle),

1) $\forall x \in S \setminus X$, $\mathrm{Omb}_X(x) = \{ s \in X \mid s \text{ lié à } x \}$ est une triade de $X$. L'application
\[
  \varphi_X : x \mapsto \mathrm{Om}_X(x) : S \setminus X \to \mathrm{Tria}(X)
\]
est une application surjective, \struck{dont les fibres sont des triades}.

2) Soit $t = \{\alpha,\beta,\gamma\}$ une triade de $X$, $\varphi_X^{-1}(\{t\})$ l'ensemble des $x \in S \setminus X$ d'ombre $t$. Alors $\bar{\bar t} = \varphi_X^{-1}(\{t\})$ est une triade \struck{notée $\varphi(t)$}, c'est l'unique triade de $S$ \add{disjointe de $t$ et} « opposée » à $t$ i.e. \uncertain{dont} les él. sont \struck{\ill{}} liés à tous les éléments de $t$ ; c'est aussi l'ens. des $s \in S \setminus t$ \struck{\ill{}} liés aux éléments de $t$. Donc $t$ est l'unique triade opposée à $\tau$.
\note{il souligne d'un trait les triades de $X$ ($t$, $t'$, $t''$) et surligne doublement l'opposée $\bar{\bar t}$ ; on ne garde que la double barre.}

3) Les deux triades $t', t''$ de $X$ \add{(liées \ill{} \ill{} par $t$)} qui sont parallèles à $t$ (sont les deux uniques triades de $S$ « liées » à $t$ \struck{disjointes} (i.e. dont les él. sont \struck{\ill{}} liés à tous les élém. de $t$)) — $t' \cup t''$ est l'ens. des $s \in S \setminus \tau$ qui sont \struck{\ill{}} liés aux él. de $\tau$.

4) Les \struck{\ill{} de $\varphi(t) \cup \varphi(t') \cup \varphi(t'')$ \ill{}} \add{triades de $X_2$}, les 6 triades $\tau, \tau', \tau''$ et $\bar\tau, \bar\tau', \bar\tau''$ d'un pavé, \ill{} les triades \uncertain{liées} forment un hexagone de triades liées

\drawing{un hexagone aux sommets $\bar\tau$ (en haut), $\tau'$, $\bar\tau''$, $\tau$ (en bas), $\bar\tau'$, $\tau''$, les côtés et les grandes diagonales en traits interrompus.}

Ceci \uncertain{indique} \ill{} parallèles les paires de triades « bi-liées » \add{« parallèles »}

\page{69}
5) Soient $\rho, \rho', \rho''$ les trois \uncertain{autres} triades de $X$, qui \uncertain{sont} \struck{mutuellement} parallèles, alors $\bar\rho, \bar\rho', \bar\rho''$ sont formées respectivement des éléments \uncertain{de l'autre} \ill{} $X$ \ill{} $\rho, \rho', \rho''$ et leur réunion, \ill{}, \uncertain{pavé} $X_2$, \uncertain{dont} \struck{elles forment} une famille de triades parallèles. Alors $X_1, X_2$ sont les seuls pavés \add{ils forment une partition de $S \setminus X$, \struck{\ill{}} (cf. card 2)} de $S$ contenus dans $S \setminus X$, \add{l'ensemble} \ill{} \uncertain{ces pavés} est en corr. biun. avec $\varepsilon$, l'ens. des deux familles de triades de $X$.

6) L'\struck{famille} \add{ens.} $I'_1$ des 3 triades de $X_1$ distinctes des $\bar\tau, \bar\tau', \bar\tau''$, et l'ens. $I'_2$ des 3 triades de $X_2$ distinctes \struck{des} \uncertain{des} $\bar\rho, \bar\rho', \bar\rho''$, est en corr. 1-1, \struck{par la} \add{les 6 triades forment un hexagone de triades liées, comme \uncertain{tout} \ill{}}, relation d'opposition. \add{Soit $I_0$ l'ens. des 3 \uncertain{parties} (= diagonales de cet hexagone)}. Pour \uncertain{tt} \add{ens. des trois \uncertain{opposées}} triades d'une des familles \add{i.e. \uncertain{tt} \uncertain{fibre} de \struck{$S \setminus X$} $\to \mathrm{Tria}(X)$} $\{\bar\tau, \bar\tau', \bar\tau''\}$ ou $\{\bar\rho, \bar\rho', \bar\rho''\}$, \ill{} bijection canonique $\sigma \simeq I_0$, \struck{\ill{}} \add{\ill{}} d'où les applications can.\ $X_1 \to I_0$ et $X_2 \to I_0$ (qui donnent lieu :
\[
  X_1 \xrightarrow{\ \sim\ } I_1 \times I_0, \qquad X_2 = I_2 \times I_0
\]
\struck{L'ens. $I_0$ est aussi l'ens. des dix triades} \add{Pour} une triade $\sigma \subset I'_1 \cup I'_2$, il f. \struck{contenues dans} \ill{} que toutes les triades de l'hexagone \ill{} des triades liées \uncertain{le} \uncertain{contiennent}, \ill{} \uncertain{seulement} le
\note{la phrase se poursuit en tête de la page 70.}

\page{70}
\ill{} \uncertain{hexagone} de triades bi-liées \uncertain{parallèles} \ill{} les \uncertain{autres} ; sont formés de triades contenues dans $S \setminus X$.

7) Considérons l'application canonique surj.
\[
  S \setminus X = X_1 \amalg X_2 \xrightarrow{\ \pi\ } I_0
\]
dont les fibres sont les réunions $\tau_1 \cup \tau_2 = \tau_1 \cup \bar\tau_1$, où $\tau_1$ est une triade de $X_1$ de famille $I'_1$, et $\tau_2$ \add{une} \struck{\ill{}} triade associée de $X_2$ \add{de famille $I'_2$} (donc $\tau_2 = \bar\tau_1$). Alors

a) Pour tout $i \in I_0$, la réunion $X \cup \pi^{-1}(\{i\})$ \add{\uncertain{est un} \ill{} contenant $X$} \struck{\ill{}} un $\mathcal{T}_1$ plongé dans $S$, et l'application $i \mapsto X \cup \pi^{-1}(\{i\})$ est une bijection de $I_0$ sur l'ensemble des $\mathcal{T}_1$ \add{$\mathfrak{X}$}. Donc les $\pi^{-1}(i)$ sont les $\mathfrak{X} \setminus X$, où $\mathfrak{X}$ est un $\mathcal{T}_1$ de $S$ contenant $X$.

8) Ces \uncertain{ens.} sont aussi les orbites dans $S \setminus X$ sous le sous-groupe de $\mathrm{Aut}(X)$ engendré par les symétries \struck{\ill{}} \add{\ill{}} \struck{triangles} de $X$. [\uncertain{Bien faire} les groupes des automorphismes « \ill{} » de $X$ \uncertain{opérant} sur $S$ …]

\section{Le diagramme des revêtements}

\page{72}
\begin{itemize}
  \item 720 : $\mathrm{Hex}$ = ens. des hexagones de triades \add{liées (de $\Delta$)} (« grands hex. ») \add{$\simeq$ ens. des pavés}
  \item 720 : $\mathrm{hex}$ = ens. des hexagones de sommets \add{de $\Delta$} \struck{\ill{}} (« petits hex. »)
  $\simeq$ ens. des paires de triades opposées
  \struck{$\simeq$ ens. des \ill{} $(X,t)$, $X$ pavé et $t \in X$}
  \item 720 : $\mathrm{tria}$ $\simeq$ \struck{ens. des triades \ill{}}
  \item 40 : $\mathrm{triHex}$ $\simeq$ ens. des triples (non ordonnés) \add{« distingués »} \struck{de} grands hexagones, i.e. \struck{\ill{}} \add{\uncertain{rayons de}} \ill{}
  \item 360 : $(\mathrm{Tria}')$ $\simeq$ ens. des \add{paires de} \struck{triades} opposées \struck{($\mathrm{Tria}' =$ ens. des \ill{})} (ce sont les diagonales des grands hexagones $\simeq$ ens. des \uncertain{faisceaux} \uncertain{réduits})
  \item 240 : \struck{\ill{}} \add{$\mathrm{Trtria}$} $\simeq$ ens. des triangles de triades parallèles (« grands triangles ») $\simeq$ ens. des \uncertain{pavés orientés}
\end{itemize}
\marginal{« 21 » en marge, en face du troisième $\simeq$ de $\mathrm{hex}$.}

\begin{tikzcd}
\mathrm{Tria}' \arrow[d, "3"'] & \mathrm{tria} \simeq \mathrm{hex} \arrow[l, "2"'] \arrow[d, "3"] \\
\mathrm{Hex} \arrow[d, "3"'] & \mathrm{Trtria} \arrow[l, "2"] \\
\mathrm{triHex} &
\end{tikzcd}

\note{dans le carré, entre les flèches verticales, la mention « (cart) » ; au-dessus de la flèche supérieure, un petit mot « \uncertain{bib} ». Le nom $\mathrm{Trtria}$ est écrit par-dessus un mot biffé.}

NB Les rev. d'ordre 2 sont les \uncertain{passages} aux triades opposées resp. aux triangles de triades \uncertain{par.} opposés.

NB Ce diagramme se déduit (\uncertain{à isom.} can. près) du diagramme
\[
  \mathrm{Tria}' \to \mathrm{Hex} \to \mathrm{triHex}
\]
\uncertain{car} le revêtement $\mathrm{Trtria}$ d'ordre 2 de $\mathrm{Hex}$ se déduit (à iso can.) du rev. d'ordre 3 $\mathrm{Hex}$ du revêtement principal de $\mathrm{triHex}$, comme le revêtement associé. D'autre part, ce diagramme
\[
  \mathrm{Tria}' \xrightarrow{\ 3\ } \mathrm{Hex} \xrightarrow{\ 3\ } \mathrm{triHex}
\]
revient \uncertain{à décrire} \struck{le \uncertain{groupe} \ill{} $I \in \mathrm{triHex}$}, $\Delta$ via la fibre $(\mathrm{Tria}')_\tau$, considérée comme un $(3,3,3)$

\page{73}
L'iso can.\ $\mathrm{hex} \simeq \mathrm{tria}$ peut s'expliciter ainsi : Soit $h$ =

\drawing{un hexagone aux sommets successifs $a$, $b'$, $c$, $a'$, $b$, $c'$.}

un hexagone, \struck{alors} formé de deux triades parallèles $t = \{a,b,c\}$, $t' = \{a',b',c'\}$. Alors $\exists !$ pavé $P \supset h$ (c'est aussi l'unique pavé contenant $\{a,b,c\}$, ou $\{a',b',c'\}$), et
\[
  t'' = P - h = \{a'', b'', c''\}
\]
est une triade par. aux deux précédentes (c'est l'unique telle). Ses él. sont en corr. 1-1 avec les diagonales de $h$, à la diagonale $\{a,a'\}$ correspondant l'él.\ $a''$ \add{\uncertain{unique}} \struck{tel que} \struck{lié} \add{non lié} à $a, a'$, \uncertain{et} lié à $b, b', c, c'$. C'est \add{(\ill{})} aussi l'unique él. de $S$ \uncertain{ayant ces propriétés} (\uncertain{c'est le troisième sommet} des triangles de côté $\{b',c\}$ — ou côté $\{b,c'\}$, — des triangles \add{dont} deux sommets \uncertain{se trouvent} \uncertain{égaux} (\uncertain{facile} : voir directement …)) et \struck{\ill{} \ill{} \ill{}} \add{\ill{} liés à $a, a'$} \struck{directement} (\uncertain{voir directement également}). Inversement, si $t'' = \{a'',b'',c''\}$ est une triade, elle est contenue dans un unique pavé $P$, et $P - t''$ est un hexagone.

Question : Soit $h$ un \add{(petit)} hexagone, correspondant à une triade $t$ (qui le détermine). Notons que $t$ détermine également un (grand) hexagone

\marginal{\ill{} \uncertain{cette} \ill{}, la troisième \ill{} \ill{} aux trois diagonales \ill{} \ill{} diagonales \ill{}}
\marginal{Autre possibilité : pour deux côtés du triangle \ill{} \ill{} ainsi 3 pts (correspondant aux \ill{} de l'hexagone) \uncertain{formant} \ill{}}
\note{ces deux notes sont écrites le long de la marge gauche, la feuille tournée.}

\page{74}
de triades $H$. Notons que le cas des hexagones \uncertain{combinatoires} \uncertain{équivaut} à : $(\mathrm{Ens})_3 \times (\mathrm{Ens})_2$, à un hexagone correspond \struck{$t$} \uncertain{l'ens.} $\tau$ de ses trois diagonales, et l'ens.\ $\varepsilon$ des deux triangles inscrits (triangles d'éléments bi-liés) — et à un couple $\tau, \varepsilon$ correspond l'hexagone $h$ dont l'ens. des sommets est $\tau \times \varepsilon$, et les sommets $(a,\alpha)$ et $(b,\beta)$ sont liés ssi $a \neq b$, $\alpha \neq \beta$. Ceci posé, l'ens.\ $\mathrm{Or}(h)$ d'un hexagone $h$ s'identifie à l'ens. des orientations du $\tau$ associé (si cette orientation est définie par $u \in$ \ill{}, celle de $h = \tau \times \varepsilon$ est donnée par $u \times u_\varepsilon$, où $u_\varepsilon$ est l'automorphisme non identique de $\varepsilon$). \add{(dans le cas d'une triade $t$ donnée)}

On explicitera (les relations entre les deux hexagones $h$ et $H$ qu'elle définit) \uncertain{en termes} d'un \ill{} du type $(3,3,3)$
\[
  I = \coprod_{\alpha \in \Lambda} I_\alpha \xrightarrow{\ p\ } \Lambda .
\]
Une triade correspond à la donnée d'un couple $i, \beta$, où $i \in I$, $\beta \in \Lambda$, $\beta \neq p(i) = \alpha$, cf. figure ; donc si $\gamma \in I$ est $\neq \alpha, \beta$, \uncertain{alors} l'ens. des trois éléments associés à $h$, qui n'est autre que $t$ lui-même, s'identifie à $I_\gamma$. Les deux triades $t', t''$ parallèles à $t$ \ill{} $\neq t$ sont évidemment celles données par $(j,\beta)$, $j \in I_\alpha \setminus \{i\}$, donc $\varepsilon \simeq I_\alpha - \{i\}$, donc
\[
  h \leftrightarrow (I_\gamma,\ I_\alpha \setminus \{i\})
\]

\drawing{trois fibres verticales de trois points chacune, marquées $\alpha$, $\gamma$ et $\beta$ (cette dernière au-dessus) ; un point de la fibre $\alpha$ est relié par des segments aux trois points de la fibre $\gamma$.}

\page{75}
D'autre part $H$ est formé de $t, t', t''$ et des triades opposées, donc
\[
  H \leftrightarrow (I_\alpha, \{0,1\})
\]
\add{« $\simeq I_\alpha \setminus \{i\}$ »} écrit sous $I_\alpha$,
\[
  \omega(H) \simeq \omega(I_\alpha) \simeq I_\alpha \setminus \{i\}, \qquad
  \omega(h) \simeq \omega(I_\gamma)
\]
Il n'y a pas de relation universelle entre ces deux ens. : 2 él.…

\uncertain{Diagramme} \add{\uncertain{Présentation}} \uncertain{symétrique} des relations entre \struck{ens.} objets de nat.\ 2, 3 de l'hexagone.

Un hexagone détermine
\begin{itemize}
  \item a) un ens.\ $\tau$ à trois éléments, qui est canoniquement l'ens. des diagonales (paires de sommets opposés) \uncertain{ou} codiagonales (paires de côtés opposés)
  \item b) l'ens.\ $\varepsilon$ (de card.\ 2) \struck{\ill{}} (l'ens. des triangles inscrits = triples de sommets bi-liés)
  \item b') l'ens.\ $\varepsilon'$ (de card.\ 2) l'ens. des triangles coinscrits = triples de côtés bi-liés.
  \item c) Un iso $\varepsilon \wedge \varepsilon' \overset{\varphi}{\simeq} \omega(\tau)$ \struck{($\simeq$ ens. des \ill{} \ill{} pour \uncertain{circ.} de $\tau$)} \struck{\ill{} $\varepsilon \wedge \varepsilon'$} (i.e. iso $\varepsilon' \simeq \omega(\tau) \wedge \varepsilon$, ou iso $\varepsilon \simeq \varepsilon' \wedge \omega(\tau)$, ou iso $\varepsilon \wedge \varepsilon' \wedge \omega(\tau) = 1_{\mathcal{T}_2}$, i.e. un él. de $\varepsilon \wedge \varepsilon' \wedge \omega(\tau)$)
\end{itemize}

On a alors
\[
  \text{Sommets} = \tau \times \varepsilon, \qquad \text{côtés} = \tau \times \varepsilon'
\]
et rel. d'incidence $(s,\alpha)$ incident à $(t,\beta)$ ssi \struck{\ill{}} $t = \varphi(\alpha \wedge \beta)$.
\note{l'indice de $1$ est la même majuscule cursive que celle transcrite $\mathcal{T}$ ; il s'agit peut-être de $\mathbb{Z}/2$. Dans la règle d'incidence, le premier terme du couple est surchargé.}

\page{76}
Le passage d'un hexagone à l'hexagone dual se traduit par le passage de $(\tau, \varepsilon, \varepsilon', \varphi)$ à
\[
  (\tau, \varepsilon', \varepsilon, -\varphi \circ \sigma)
\]
où $\sigma$ est la symétrie $\varepsilon' \wedge \varepsilon \to \varepsilon \wedge \varepsilon'$.
\[
  \mathrm{Aut}(h) \subset \mathrm{Aut}(\tau) \times (\pm 1)^2
\]
\add{$(\pm1)^2 = \mathrm{Aut}\,\varepsilon \times \mathrm{Aut}\,\varepsilon'$} formé des triples $u, \alpha, \beta$ tels que
\[
  \alpha\beta = \mathrm{sg}(u) \quad\text{ou}\quad \alpha\beta\,\mathrm{sg}(u) = 1
\]
i.e. c'est le sous-groupe de $\mathrm{Aut}\,\tau \times (\pm 1)^2$ formé \uncertain{image} \uncertain{inverse} de
\[
  (\pm 1)^{3\prime} = \{ (\alpha,\beta,\gamma) \in (\pm 1)^3 \mid \alpha\beta\gamma = 1 \}.
\]
\note{l'argument de $\mathrm{sg}$ est écrit comme un $\tau$ ; on lit $u$ d'après les triples $u, \alpha, \beta$. Les lettres $\alpha$ et $\beta$ du premier membre sont surchargées.}

\drawing{un cube de revêtements, flèches marquées de leur degré (2 ou 3). Sommets lisibles : « $\tau$ = diagonales » relié par 2 à « sommets » ; « côtés » ; « rep.\ sym. » ; « sommets orientés $\simeq$ côtés $\simeq$ repères » ; $\varepsilon$ ; « triangles inscrits » ; « $\varepsilon' = \varepsilon \wedge \omega$, triangles coinscrits » ; « $\varepsilon \times \omega$, triangles inscrits orientés » ; « $\omega = \omega(\tau)$, orientations » ; le point $e$. Les flèches convergent vers $e$ ; plusieurs flèches et étiquettes du coin supérieur gauche sont surchargées et illisibles.}

NB Tous les carrés sont cartésiens, y compris les petits carrés horizontaux de forme triangulaire $\nwarrow \leftarrow \uparrow$.

\page{77}
\note{feuille tenue en largeur et très chargée : à gauche, un dessin au crayon d'un hexagone avec ses diagonales, un schéma reliant $t_1, t_2, t_0$ à $\alpha, \beta, \gamma$ et le tableau de couples $(t_1,\alpha)(t_1,\beta)(t_1,\gamma)$, $(t_2,\alpha)(t_2,\beta)(t_2,\gamma)$, $(t_0,\alpha)(t_0,\beta)(t_0,\gamma)$, surmonté de $p_1 \dots p_6$ ; on transcrit la liste encadrée à droite, puis on décrit le grand diagramme.}

$H$ un hexagone

\begin{itemize}
  \item $H_d$ ens. des diagonales $\simeq$ ens. des codiagonales — 3
  \item $H_t$ ens. des triangles inscrits — 2
  \item $H_\omega$ ens. des orientations $\simeq \omega(H_d)$ — 2
  \item \add{$H_{t^*}$ ens. des triangles coinscrits} — 2
  \item $H_s$ ens. des sommets $= H_d \times H_t$ — 6
  \item $H_{s^*}$ ens. des côtés $= H_d \times H_{t^*}$ — 6
  \item $H_r$ ens. des repères $\subset$ \uncertain{$H_s \times H_d$} (« rel. d'inc. ») — 12
  \item \struck{$H_{r'}$ ens. des \ill{} $=$} $H_d \times H_d^* \times H_t \times H_{t^*}$ — \struck{6}
  \item $H_{tro}$ = triangles inscrits orientés
\end{itemize}
\note{en face de $H_\omega$, une formule entièrement surchargée finissant par $\omega(H_d)$.}

\drawing{grand diagramme de revêtements, flèches marquées 2 ou 3, entre : « hex + diag $\simeq$ tria + él. marqué $\simeq$ couples de triades \uncertain{séquentes} (\uncertain{explicites}) » ; « $\mathrm{tria}' = \mathrm{Hex} + \mathrm{diag}$, $\mathrm{Hex} + \mathrm{codiag} = \mathrm{bib.\,plié}$ » ; « $\mathrm{tria} \simeq \mathrm{hex} = \mathrm{Hex} + \text{sommet}$ » (flèche « opp ») ; « hex or (hex + orientation) $\simeq$ tria or (triade + ordre circulaire) $\simeq \omega((\mathrm{hex+diag})/\mathrm{hex})$ » ; « $\mathrm{Hex}$ » ; « $\mathrm{trtria} \simeq \omega(\mathrm{Hex}/\mathrm{triHex}) \simeq$ Hex + triangle inscrit $\simeq$ pavé orienté » ; « hex + sommet » ; « hex + triangle inscrit orienté $\simeq$ couple de triades orientées parallèles » ; « $\mathrm{Hex} + \mathrm{rep} \simeq$ couple de triades liées $\simeq$ couple de triades par. $\simeq$ hex + triade marquée $\simeq$ hex + triangle inscrit » avec les mentions « torseur sous Hex de groupe $\mathfrak{S}_3 \times \mathfrak{S}_2$ », « torseur sur trtria de groupe $\mathfrak{S}_3$ », « torseur sur tria de groupe $\mathfrak{S}_2$ » ; « $\mathrm{tria}' = \mathrm{Hex} + \mathrm{rep\ mod\ opp}$ (torseur sur Hex de groupe $\mathfrak{S}_3$) » ; « $\mathrm{triHex}$ », « $\mathrm{triHex}$ » (doublé) ; « $\mathrm{Hexor} = \omega(\mathrm{tri}'/\mathrm{Hex}) = \mathrm{Hex} + \mathrm{or}$ » ; « dror tria $\simeq$ Hex + dr.\ or.\ \uncertain{inscr.} ». En bas, un cube des ensembles $H_s$, $H_r$, $H_{p'}$, $H_{s^*}$, $H_d$, $H_{tr}$, $H_{trr}$, $H_\omega$, $H_{tr^*}$, $e$, marqué $\mathfrak{S}_3 \times \mathfrak{S}_2$, avec de petits hexagones dessinés à côté.}

$\mathfrak{S}_3 \times \mathfrak{S}_2$ :
\begin{itemize}
  \item \struck{7} \add{3} s.-groupes d'indice 2 (\uncertain{3 classes})
  \item 3 s.-groupes d'indice 3 (\uncertain{conj.})
  \item 3 sous-groupes d'indice 6 \add{\uncertain{non}} (\uncertain{conjugués})
\end{itemize}

\page{78}
\section{Étude combinatoire des polygones}

Étude combinatoire des polygones de $3 \leq n \leq 6$ côtés

1) $n=3$ \quad $\mathrm{Pol}_3 \simeq (\mathrm{Ens})_3$, le foncteur dualité est isom. can. à l'identité (si $P \in \mathrm{Ob}\,\mathrm{Pol}_3$, alors $P^\circ \simeq P$ \uncertain{iso can.}) car les côtés sont en corr. 1-1 avec les sommets. (Ceci se généralise aux polygones à un nb impair de côtés. Groupe : $\mathfrak{S}_3$)

2) $n = 4$ \quad $\mathrm{Pol}_4 \simeq (\mathrm{Ens})_{2,2}$ (se généralise aux \uncertain{cubes} de dim $d \in \mathbb{N}$,
\[
  \mathrm{Cube}_d \simeq (\mathrm{Ens})_{\underbrace{2,2,\dots,2}_{d}} \Big)
\]
\struck{\ill{}} \uncertain{Associer} au \uncertain{carré} : \uncertain{tt} carré comb.\ $C$ l'ens. de ses quatre côtés, avec la division en côtés opposés. Le \uncertain{cube} combinatoire \uncertain{est décrit} par l'ens. des 4 sommets, avec division en classes de sommets \struck{opposés}. Or les sommets correspondent aux sections de $I$ sur $\Lambda$, l'incidence signifiant que les \struck{côtés} sections ne se rencontrent pas, l'\uncertain{opposition} \uncertain{ou} \uncertain{non-incidence} signifiant « parallèles »…

\drawing{deux fibres $I_\alpha$, $I_\beta$ de deux points chacune, au-dessus de $\Lambda$.}

\begin{tikzcd}[column sep=small, row sep=small, nodes={font=\scriptsize}]
C_s \arrow[d, "2"'] & & C_r \arrow[ll, "2"'] \arrow[dl, "2"] \arrow[d, "2"] \\
C_d \arrow[dd, "2"'] & C_{p'} \arrow[l, "2"'] \arrow[d, "2"] \arrow[ddr, "2"] & C_{s^*} \arrow[dd, "2"] \\
 & C_\omega \arrow[dl, "2"'] & \\
e & & C_{d^*} \arrow[ll, "2"]
\end{tikzcd}

\note{diagramme redessiné d'après un croquis serré, où il est tracé en perspective : les flèches obliques portent les mentions « opp » et « cart » (carrés cartésiens) ; la redisposition sur une grille garde les sommets et les degrés, mais l'attribution de quelques flèches autour de $C_\omega$ et $C_{d^*}$ est douteuse.}

\begin{itemize}
  \item $C_r$ ens. des repères
  \item $C_s$ — sommets
  \item $C_{s^*}$ — côtés
  \item $C_d$ — diagonales $= C_s/\mathrm{opp}$
  \item $C_{d^*}$ — codiagonales $= C_{s^*}/\mathrm{opp}$
  \item $C_{p'}$ — repères \uncertain{mod} symétrie $= C_r/\mathrm{opp}$
  \item $C_\omega$ — orientations
\end{itemize}

$C_r$ est un torseur sous $D_4 \simeq \mathfrak{S}_2 \cdot \mathbb{Z}/4\mathbb{Z}$. On trouve que si $r \in C_r$ et si on \uncertain{prend} les stabilisateurs des images de $r$ dans les ens. du diagramme, on trouve exactement \ill{} fois \ill{} une seule \uncertain{classe} \ill{} sous-groupe de $D_4$ [il y a \uncertain{exactement} 8 \uncertain{classes}, y compris \ill{} et $D_4$

\marginal{On obtient les \uncertain{sous-groupes} comme stabilisateurs (\ill{} les côtés) de chaque él. de \ill{} $\omega(T)$}
\marginal{\ill{} que $\mathbb{Z}$ \ill{} 4 d'indice 2 qui \ill{} \ill{} conjugués, engendrés par \ill{} $(\varepsilon.\dot1, \varepsilon.\dot3)$, $(\varepsilon.\dot{?}, \varepsilon.\dot2)$…}

\page{79}
3) $n = 6$, \quad $P_6 \simeq P_3 \times (\mathrm{Ens})_2$ (\uncertain{marche} chaque fois qu'on a $P_n$ avec $n = 2m$, $m$ impair…) d'où
\[
  P_6 \simeq (\mathrm{Ens})_3 \times (\mathrm{Ens})_2
\]
Si $t \in (\mathrm{Ens})_3$ et $\omega \in (\mathrm{Ens})_2$, \struck{\ill{}} l'hexagone associé a comme sommets les él. de $t \times \omega$, et $(a,\alpha)$, $(b,\beta)$ \struck{\ill{}} liés ssi $a \neq b$, $\alpha \neq \beta$ : l'hexagone est une bitriade. \uncertain{Inv.}, pour un hexagone donné $H$, \add{si $\mathrm{opp}$ est la symétrie antipodale} il y a bijection canonique entre l'ens.\ $H_s/\mathrm{opp}$ \struck{\ill{}} des diag. et l'ens.\ $H_{d^*}$ des \struck{\ill{}} codiag. [paires de côtés opposés] en associant à chaque diagonale la paire des 2 côtés parallèles à celle-ci.
\[
  \boxed{H_d \simeq H_{d^*}} \quad (= t)
\]
Soit $H_{tr}$ l'ens. des 2 « triangles inscrits » — c'est $\varepsilon$ — \add{ou l'ens. des triades constituant cet hexagone} — \ill{} et $H_{tr^*}$ l'ens. des triangles coinscrits.

On a maintenant le diagramme

\drawing{un cube de revêtements analogue à celui de la page 78 : $H_r$ (marqué $D_6$) $\xrightarrow{2} H_s$, $H_r \to H_{p'}$ (marqué $\mathfrak{S}_3$), $H_{s^*}$, $H_d$ (flèches « opp » et 2), $H_{trr}$, $H_{tr}$ (marqué $\mathbb{F}_2^2$), $H_\omega$, $H_{tr^*}$, $e$ ; flèches verticales de degré 3, les autres de degré 2. À gauche, un petit hexagone à triangles inscrits en pointillés.}

NB Dans les carrés, \uncertain{contenant} les flèches hor. de degré 2 \ill{} 3

NB
\begin{itemize}
  \item $H_d$ = ens. des \uncertain{sommets} d'un triangle inscrit \uncertain{ou} coinscrit
  \item $H_\omega$ = ens. des orientations de $H \simeq \omega(H_d) \simeq H_{tr} \wedge H_{tr^*}$
  \item $H_{p'} = H_r/\mathrm{opp} \simeq$ ens. des repères de $H_d$ (dans $\mathfrak{S}_3$-torseurs)
  \item $H_{tr} \times H_{tr^*} \simeq H_{p'} \times H_\omega \simeq H_{trr}$ ens. des 4 triangles inscrits \uncertain{orientés}
\end{itemize}
Ceci correspond à un diagramme de sous-groupes de
\[
  D_6 \simeq \mathfrak{S}_2 \cdot \mathbb{Z}/6\mathbb{Z} \simeq \mathfrak{S}_2 \times \mathfrak{S}_3
\]
d'ordre 12. $D_6$ \uncertain{ens.} + \add{trois} \struck{\ill{}} sous-groupes \add{(conjugués)} d'indice 3 (stabilisateurs des \add{trois} él. de $H_d$), 3 sous-groupes d'indice 2 \add{(\uncertain{non} conjugués)} (stabilisateurs des él. de $H_{tr}$, $H_{tr^*}$, $H_\omega$ respectivement), un sous-groupe d'indice 4 (stabilisateur des él. de $H_{trr}$)

\page{80}
\uncertain{enfin} trois \struck{familles} \add{classes de conjugaison} (de sous-groupes \add{d'ordre 2}, \struck{\ill{}} \add{6} \uncertain{comprenant} un groupe invariant \uncertain{qui} forme une seule classe), savoir le \uncertain{centre}, qui est le stabilisateur des él. de $H_{p'}$, et 2 \struck{familles} \add{classes} de \add{3} s.-groupes chacune \struck{\ill{}} resp. stabilisateurs des él. de $H_s$ et $H_{s^*}$ : les \struck{\ill{}} \add{et les 2n} (\uncertain{12} él. \ill{} tour des côtés, i.e. gén. du facteur $\mathbb{Z}/2\mathbb{Z} = \mathfrak{S}_2$ de $\mathfrak{S}_2 \times \mathfrak{S}_3$, \uncertain{ou} él. d'ordre 2 de $\mathfrak{S}_3$) \struck{\ill{}}

Les classes de conj. de s.-groupes de $D_6$ correspondent \add{1-1} \uncertain{aux} \struck{\ill{}} \uncertain{espaces} \uncertain{homogènes} \ill{} du diagramme \uncertain{marqué} ; et on constate que les \uncertain{automorphismes} de ces espaces homogènes sont \uncertain{aussi} marqués sur le diagramme (une flèche $\xrightarrow{2}$ i.e. de degré 3 \uncertain{impliquant} qu'il y a 2 \uncertain{non} \ill{} automorphismes \ill{}) …
\note{la page s'arrête là ; la suite, si elle existe, est hors de ce lot.}

\end{document}
